\chapter{Vision Screening}
\section*{What Is This Test?} Vision screening uses brief tests to identify reduced visual acuity, eye misalignment, refractive error, or other concern needing eye examination.\section*{Alternative Names} Visual acuity screening; eye screening; amblyopia screening.\section*{Principle} Charts, photoscreeners, autorefractors, and alignment tests compare visual function with age-appropriate expectations.\section*{Why This Test Is Ordered} It is used in children and adults to identify refractive error, amblyopia risk, eye disease, or vision change.\section*{Primary vs Secondary Test} Screening is preliminary; comprehensive examination by an eye professional provides diagnosis.\section*{How This Test Is Performed} Each eye is tested separately and together for acuity, alignment, color or depth when indicated, and sometimes with an instrument.\section*{What Preparation Is Needed} Bring glasses or contact-lens information; no fasting is required.\section*{Understanding Results} A failed screen means further examination is needed, not necessarily disease. Normal screening does not exclude every eye disorder.\section*{Follow-up Tests} Refraction, dilated examination, pressure measurement, visual fields, OCT, or ophthalmology referral may follow.\section*{References} \begin{enumerate}\item MedlinePlus. Vision Screening.\item American Academy of Ophthalmology. Vision screening recommendations.\end{enumerate}\clearpage\section*{Understanding Results and Follow-up}
The laboratory report should identify the specimen, method, units, reference interval or decision limit, and whether the result is positive, negative, abnormal, or inconclusive. Reference intervals vary with age, sex, pregnancy, specimen type, assay, and laboratory; a result outside a range is not automatically a diagnosis, and a result within a range does not always exclude disease. Discuss unexpected results with the ordering clinician, who can decide whether repeat or confirmatory testing, treatment, or referral is appropriate. Do not change medicines or delay urgent care based on an isolated result.
\section*{Regulatory Status and Authorization Date}This chapter describes a test category, not a specific commercial product. FDA, CDSCO, EU IVDR, and MHRA approval or registration dates therefore cannot be assigned without the assay manufacturer, product name, instrument, and intended use. For laboratory-developed tests, an individual agency approval date may not exist. See the Preface for the interpretation of regulatory status.
\clearpage
